\chapter{Measurements}
\label{ch:measured}

Three new measurements are reported here: how much of a transition each form of label keeps, how
the state machine behaves when a fixed scenario is put to a real host many times, and whether the
side a branch is read from leaves a mark. The transition tables are read directly from the
language's definition for each measurement, never transcribed.

\section{What a label keeps of its transition}
\label{sec:labels}

\textbf{Method.} A transition is its pair, the state left and the state reached. A branch that
reads only the label recovers the pair exactly where the label is reached from one pair, and loses
information where it is shared. The loss is the conditional entropy
\(H(\text{pair} \mid \text{label})\) \cite{src:Shannon}, and the label keeps
\(H(\text{pair}) - H(\text{pair} \mid \text{label})\), the mutual information between the two.
The 63 named cells of Table~\ref{tab:matrix} are weighted alike, \(H(\text{pair}) = 5.977\) bits.
Each pair is put through every situation a label is decided in, four kinds of answer at four
costs for each side, against a timeout of 1000 and a held edge of 120 units, with every situation
weighted alike within its pair.

\textbf{Results.} Table~\ref{tab:labelbits} gives four forms of label.

\begin{table}[h]
\centering
\small
\begin{tabular}{@{}l r r r@{}}
\toprule
form of label & labels & \(H(\text{pair}\mid\text{label})\) & kept \\
\midrule
one name a cell, as the tables print it & 19 & 2.675 & 3.302 \\
one name, decided in its situation & 26 & 2.129 & 3.848 \\
sync written as its passage & 41 & 1.121 & 4.856 \\
gray trips also carrying their state & 53 & 0.498 & 5.480 \\
\bottomrule
\end{tabular}
\caption{Bits of a transition each form of label keeps, of 5.977.}
\label{tab:labelbits}
\end{table}

Read from the tables alone, sync is reached from 22 pairs and leaves 4.459 bits of its pair
unread; fizz and fuzz, 7 pairs each, leave 2.807; drop, 6 pairs, 2.585; halt, 5 pairs, 2.322; join
and wake, 2 pairs each, 1.000.

Deciding each label in its situation moves six of sync's pairs to the high load table's names,
surg, vent twice, snap, lead and rite, and leaves sync reached from 16 pairs at 3.999 bits.
Writing sync as its passage, fuzz carrying the state left and fizz the state reached, removes its
loss entirely.

What remains, 1.121 bits, sits in eight labels: fizz and fuzz into and out of gray, 2.807 bits
each; drop, 2.539; halt, 2.322; join, tilt, vent and wake, 1.000 each. A gray trip loses what a
sync passage keeps because its labels carry no state. Letting fuzz into gray carry the state it
left, and fizz out of gray the state it reaches, raises what a label keeps to 5.480 bits of
5.977. The loss left, 0.498 bits, sits in drop, halt, join, tilt, vent and wake. drop and halt
are the largest: each is a rule, any stable state to void and any state to blok, and a rule
covering many pairs names none of them.

\section{A scenario put to a real host}
\label{sec:scenario}

\textbf{Method.} A scenario fixes, for each of ten cycles, what each side is: held, not held,
late, or not asked. Each side is realized as a run of a million asks between two readings of the
host clock, after four baseline cycles from which the timeout and the busy edge are derived. The
host is a Windows workstation.

\textbf{Result.} The clock printed by one run is
\begin{flushleft}\small\ttfamily
[dual] -(drop)-> [lead] -(pass)-> [rite] -(join)-> [dual] -(fuzz+dual)-> [sync] -(fizz+wait)->
[wait] -(loss)-> [void] -(fuzz)-> [gray] -(fizz)-> [dual] -(fuse)-> [void] -(sprk)-> [dual]
\end{flushleft}
In the field's terms: both sides healthy, the right side lost, failover to the right, both healthy
again, the left side lagging, the lagging side timing out into an outage, the left side going
unobserved and returning, the breaker tripping, and recovery. Every state agrees with what the
recorded sides imply.

\section{The scenario repeated}
\label{sec:repeat}

\textbf{Method.} The scenario is put 40 times in succession, about 20 seconds in all, and every
trace is read back. Two batches of 40, A and B, were put minutes apart.

\textbf{Results.} Table~\ref{tab:repeat}.

\begin{table}[h]
\centering
\small
\begin{tabular}{@{}l r r@{}}
\toprule
reading & A & B \\
\midrule
cycles read as other than the recorded sides imply & 0 of 400 & 0 of 400 \\
held asks past the timeout their run derived & 2 of 480 & 0 of 480 \\
both-held cycles read busy in place of dual & 8 of 160 & 15 of 160 \\
distinct clocks & 5 & 8 \\
runs printing the most common clock & 32 of 40 & 28 of 40 \\
derived timeout, least & 72,182 & 79,782 \\
derived timeout, median & 150,810 & 167,807 \\
derived timeout, most & 191,351 & 212,558 \\
\bottomrule
\end{tabular}
\caption{Two batches of forty runs of the scenario. Timeouts are in counts of the host clock.}
\label{tab:repeat}
\end{table}

\textbf{The machine reads its record exactly.} In all 800 cycles the state is one the recorded
sides imply. Every departure from the scenario is the host's, in timing.

\textbf{Spurious timeouts.} Two held asks of 960, 0.21\%, came in past the timeout their run
derived: correct answers, late. Both were in batch A, both in the cycle where a held side stands
beside an unasked one and the pair reads gray whatever the held side does, and both in runs whose
derived timeouts were among the lowest, 72,910 and 74,596 against a median of 150,810.

\textbf{The timeout varies.} Within a batch the largest derived timeout is 2.65 times the least in
A and 2.66 times in B. Each run derives its own from the host as it is at that moment, and the
spread across runs measures how far a figure written once would have been wrong.

\textbf{busy varies; nothing else does.} Of 320 both-held cycles, 23 read busy, 5\% in A and 9.4\%
in B. The scenario does not choose busy; the host's load does. The distinct clocks differ in these
readings only. In A each of the four both-held cycles reads busy in 2 runs and dual in 38; in B in
3, 2, 6 and 4 runs. Every other cycle reads the same state in all 80 runs: lead, rite, void, wait
and gray, each decided by which sides held, never varied.

\textbf{busy reads the host, not a side.} In batch B, 53 of the 320 sides of the both-held cycles
came in past the edge, a share of 0.1656. Were the two sides independent, both would be past it at
once in \(160 \times 0.1656^2 = 4.39\) cycles. They were past it together in 15, 3.4 times as often.

\section{Which side is asked first}
\label{sec:side}

\textbf{Method.} core and surv, lead held and rite held, are distinct only where the side a branch
is read from leaves a mark. Two branches doing the same work are put through the known order of
Proposition~\ref{prop:slice}, the branch asked first alternating from trial to trial, and every
pass is solved in exact integers.

\textbf{Results.} On the host the branch asked first reads dearer in 12 of 29 trials that are not
even, with no lean past twice the spread. With one branch made 512 reads dearer, the cheaper reads
cheaper in 46 to 58 of 64 comparisons at every link, from both sides alike. On this host a steady
lead and a steady rite are one state seen from two sides; on a part whose side does leave a mark
they are two.

\section{Checks of the definition}

On traces whose answers are known, every rule of Chapter~\ref{ch:states} holds. Every pair the
tables name takes a label from its own candidates in all 256 situations; gray trips are atomic
over 500 random traces of 40 cycles; and all 3,035 sync passages in those traces read back their
exact pair from their labels. Of the 49 transitions among the first seven states, 12 take a name
nothing else uses and sync covers 22. Where both tables name a transition, 3 agree, the second
names 7 the first calls sync, and 2 disagree.
